% =====================================================================
%  Курсовая работа «Разработка межсетевого экрана»
%  Задача 2. Архитектура ПО МЭ — описание к схеме architecture_firewall.drawio
%
%  Компиляция: pdfLaTeX, дважды.
%  Рисунок: в draw.io «Файл -> Экспорт -> PDF» (галочка «Обрезать»),
%  сохранить как architecture.pdf рядом с этим файлом.
% =====================================================================
\documentclass[12pt,a4paper,final]{article}

\usepackage[T2A]{fontenc}
\usepackage[utf8]{inputenc}
\usepackage[english, russian]{babel}
\usepackage{amsmath}
\usepackage[left=3.00cm, right=2.00cm,top=3.00cm,bottom=3.00cm]{geometry}
\usepackage{graphicx}
\usepackage{setspace}
\onehalfspacing
\usepackage{float}
\usepackage{pdflscape}
\usepackage{booktabs}
\usepackage{array}
\usepackage{longtable}
\usepackage{enumitem}
\usepackage{caption}
\captionsetup{labelsep=endash,justification=centering}
\captionsetup[longtable]{labelsep=endash,justification=raggedright,singlelinecheck=false}
\renewcommand{\arraystretch}{1.15}
\newcolumntype{C}[1]{>{\centering\arraybackslash}p{#1}}
\newcolumntype{L}[1]{>{\raggedright\arraybackslash}p{#1}}
\setlength{\emergencystretch}{3em}
\tolerance=2000

\usepackage[unicode, pdftex, hidelinks]{hyperref}

\begin{document}
	\noindent
	\begin{center}
		\thispagestyle{empty}
		\large{МИНИСТЕРСТВО НАУКИ И ВЫСШЕГО ОБРАЗОВАНИЯ РОССИИ\\
			\vspace{3mm}
			ФЕДЕРАЛЬНОЕ ГОСУДАРСТВЕННОЕ БЮДЖЕТНОЕ ОБРАЗОВАТЕЛЬНОЕ}\end{center}\vspace{-5mm}
	\begin{center}\large{УЧРЕЖДЕНИЕ}\end{center}
	\begin{center}\vspace{-3mm}\large{ВЫСШЕГО ПРОФЕССИОНАЛЬНОГО ОБРАЗОВАНИЯ}\end{center}
	\begin{center}\large\textbf{<<САНКТ-ПЕТЕРБУРГСКИЙ ПОЛИТЕХНИЧЕСКИЙ УНИВЕРСИТЕТ
			ПЕТРА ВЕЛИКОГО>>}\end{center}
	\begin{center}
		\normalsize{Институт Компьютерных наук и Кибербезопасности}
	\end{center}
	\begin{center}
		Высшая школа технологий искусственного интеллекта
	\end{center}
	\begin{center}
		Направление 02.03.01 Математика и компьютерные науки
	\end{center}

	\vspace{0.3cm}

	\begin{center}
		\large Защита информации
	\end{center}
	\begin{center}
		\Large\textbf{Архитектура программного обеспечения}
	\end{center}
	\begin{center}
		\large<<Разработка межсетевого экрана>>
	\end{center}

	\begin{flushleft}
		Студенты группы
	\end{flushleft}
	\vspace{-0.55cm}
	\begin{flushleft}
		\hbox to 5.5cm {№5130201/30102:\quad \hrulefill}
	\end{flushleft}
	\begin{flushright}
		\vspace{-1.45cm}
		Смирнов С.Р.
	\end{flushright}
	\begin{flushright}

		Михайлова А.А.
	\end{flushright}
	\begin{flushright}

		Аббад М.С.
	\end{flushright}
	\vspace{1cm}
	\begin{flushleft}
		\hbox to 5.5cm {Преподаватель:\quad \hrulefill}
	\end{flushleft}
	\begin{flushright}
		\vspace{-1.5cm}
		Силиненко А.В.
	\end{flushright}
	\begin{flushright}
		\vspace{1cm}
		\hspace{5cm}\hbox to 1.2cm {<<\hrulefill>>}\hbox to 2.2cm {\hrulefill }\hbox to 1.5cm { 20\hrulefill г.}
	\end{flushright}
	\begin{center}
		\vspace{4cm}
		Санкт-Петербург 2026
	\end{center}
	\newpage

\tableofcontents
\clearpage

% =====================================================================
\section{Общая структура}\label{sec:overview}

ПО МЭ состоит из четырёх подсистем, которые выделены в ОБЩ-04: фильтрации,
управления, регистрации и контроля целостности. Им служат два вспомогательных
слоя: постоянные данные и среда ОС Linux. Схема приведена на
рисунке~\ref{fig:arch}, назначение частей кратко описано в
таблице~\ref{tab:parts}.

\begin{landscape}
\begin{figure}[p]
\centering
\IfFileExists{architecture.pdf}{%
  \includegraphics[width=\linewidth,height=0.86\textheight,keepaspectratio]{architecture.pdf}}{%
\IfFileExists{architecture.png}{%
  \includegraphics[width=\linewidth,height=0.86\textheight,keepaspectratio]{architecture.png}}{%
  \fbox{\parbox{0.8\linewidth}{\centering\vspace{2cm}
  Экспортируйте схему из draw.io в файл \texttt{architecture.pdf}\\
  и положите его рядом с этим документом.\vspace{2cm}}}}}
\caption{Архитектура ПО МЭ}\label{fig:arch}
\end{figure}
\end{landscape}

{\small
\begin{longtable}{|L{3.6cm}|L{11.0cm}|}
\caption{Части ПО МЭ}\label{tab:parts}\\
\hline
\textbf{Часть} & \textbf{Назначение}\\
\hline\endfirsthead
\hline
\textbf{Часть} & \textbf{Назначение}\\
\hline\endhead
01 Подсистема фильтрации & Принимает кадры с портов, ищет сеанс в таблице
  состояний, выполняет DNAT, проверяет политику доступа, выполняет SNAT и передаёт
  кадр в выходной порт. Работает непрерывно и не зависит от того, работает ли
  управление.\\ \hline
02 Подсистема управления & Три точки входа администратора (CLI на консоли,
  CLI по SSH, web-интерфейс по HTTPS) и общее ядро: единый API управления
  и прикладные сервисы. Только через неё администратор меняет состояние ПО МЭ.\\ \hline
03 Подсистема регистрации & Принимает события от подсистем и записи о пакетах
  от фильтрации, хранит журналы, выполняет выборки.\\ \hline
04 Подсистема контроля целостности & Считает SHA-256 контролируемых файлов,
  сравнивает с эталонами и при нарушении включает глобальную блокировку трафика.\\ \hline
Постоянные данные & Конфигурация, учётные записи, эталонные контрольные
  суммы с резервными копиями, два журнала. Всё хранится на диске и переживает
  перезагрузку.\\ \hline
ОС Linux / жизненный цикл & Интерфейсы, IP-адреса, маршруты, время, запуск
  и перезапуск процессов, выключение и перезагрузка изделия.\\ \hline
\end{longtable}
}

Архитектура построена на трёх принципах.
\begin{enumerate}[label=\arabic*)]
\item \textbf{Каждая подсистема~— отдельный процесс.} Подсистему можно
  остановить и перезапустить, не трогая остальные (ОБЩ-04, НДЖ-01). Ошибка
  в web-интерфейсе не останавливает фильтрацию.
\item \textbf{Одна логика для CLI и web.} Точки входа только принимают команды
  и показывают результат. Проверка прав, валидация и выполнение сосредоточены
  в едином API управления, поэтому функции в CLI и web совпадают (УПР-02),
  а правила проверяются одинаково.
\item \textbf{Отказ в «закрытое» состояние.} Если фильтрация не работает или
  включена глобальная блокировка, трафик через изделие не проходит (НДЖ-02,
  ЦЛС-04).
\end{enumerate}

% =====================================================================
\section{Подсистема фильтрации}\label{sec:filter}

\subsection{Тракт пакета}

Кадр проходит через шесть блоков слева направо (верхняя часть
рисунка~\ref{fig:arch}). Блоки перечислены в таблице~\ref{tab:pipeline}.

{\small
\begin{longtable}{|C{0.6cm}|L{3.1cm}|L{8.3cm}|L{2.3cm}|}
\caption{Блоки тракта фильтрации}\label{tab:pipeline}\\
\hline
\textbf{№} & \textbf{Блок} & \textbf{Что делает} & \textbf{Требования}\\
\hline\endfirsthead
\hline
\textbf{№} & \textbf{Блок} & \textbf{Что делает} & \textbf{Требования}\\
\hline\endhead
1 & Приём / передача Ethernet &
  Получает кадр с порта \texttt{port1}\,\dots\,\texttt{port6} через сокет
  \texttt{AF\_PACKET}, ставит метку времени, разбирает заголовки, собирает
  фрагменты. При включённой глобальной блокировке кадр сразу удаляется.
  Тот же блок отвечает за сопоставление логических имён портов именам
  интерфейсов ОС. &
  ОБР-01, ОБР-03, ОБР-06, ОБР-07, ЦЛС-04\\ \hline
2 & Контроль состояний &
  Ищет сеанс в таблице состояний (TCP, UDP, ICMP) в прямом и обратном
  направлении. Если сеанс известен, кадр транслируется по записи и сразу
  уходит в выходной интерфейс. Если нет~— продолжает путь как новый. &
  СОС-02, СОС-03\\ \hline
3 & DNAT (до политики) &
  Проверяет таблицу DNAT; первое подходящее правило заменяет адрес и порт
  получателя. &
  ФЛТ-10, НАТ-01\,--\,НАТ-04\\ \hline
4 & Движок политики &
  Просматривает правила по возрастанию номеров и применяет первое подходящее
  активное правило: «передать», «перейти» или «удалить». Политика видит адрес
  получателя \emph{после} DNAT. &
  ФЛТ-01\,--\,ФЛТ-09\\ \hline
5 & SNAT (после политики) &
  Для пропущенного кадра проверяет таблицу SNAT и заменяет адрес и порт
  отправителя. Если у правила политики включён контроль состояния, создаёт
  запись нового сеанса с исходными и транслированными адресами. &
  ФЛТ-10, ФЛТ-11, СОС-01\\ \hline
6 & Выходной интерфейс &
  Пересчитывает контрольные суммы, определяет MAC-адрес следующего узла
  (ARP, маршруты) и передаёт кадр в назначенные порты или в \texttt{local}
  (TAP-интерфейс к собственным службам изделия). Если порт выключен или
  не назначен, кадр удаляется. &
  ФЛТ-08, ОБР-05, ОБР-08\\ \hline
\end{longtable}
}

\subsection{Почему DNAT стоит до политики, а SNAT после}

Такой порядок задан в НАТ-02: сначала DNAT, затем политика с учётом
изменённого адреса, затем SNAT. Смысл его в следующем.

\textbf{DNAT до политики.} Пусть извне приходит запрос на адрес изделия
\texttt{203.0.113.5:8080}, а правило DNAT перенаправляет его на внутренний
сервер \texttt{192.168.1.10:80}. Поскольку DNAT выполнен раньше, политика
видит реальный адрес получателя \texttt{192.168.1.10:80}. Администратор
пишет правило «разрешить TCP к \texttt{192.168.1.10} порт 80», то есть
описывает доступ к реальному серверу, а не к адресу, который в сети
нигде не существует. Кроме того, от адреса получателя после DNAT зависит
выбор выходного порта.

\textbf{SNAT после политики.} SNAT заменяет адрес отправителя, например
внутренние адреса \texttt{192.168.1.0/24} на внешний адрес изделия. Если
выполнять его после политики, то политика видит исходного отправителя
и может разрешать или запрещать доступ по внутренним подсетям и
конкретным узлам. Кроме того, для SNAT с адресом выходного интерфейса нужно
знать этот интерфейс, а его выбирает только политика (действие
«передать» с указанием выходных портов).

В результате каждый вид трансляции выполняется ровно там, где для него есть
вся нужная информация, а политика всегда работает с «настоящими» адресами
узлов.

\subsection{Уже известный сеанс}

Верхняя пунктирная стрелка показывает быстрый путь. Для пакетов уже
установленного сеанса DNAT, политика и SNAT повторно не выполняются: запись
в таблице состояний уже хранит обе пары адресов, до и после трансляции.
\begin{itemize}
\item Пакеты в прямом направлении получают ту же трансляцию, что и первый
  пакет сеанса.
\item Ответные пакеты получают обратную трансляцию: например, ответ сервера
  \texttt{192.168.1.10:80} возвращается клиенту с адресом отправителя
  \texttt{203.0.113.5:8080}.
\end{itemize}

Нижняя пунктирная стрелка (от SNAT к контролю состояний) показывает, когда
запись создаётся: \emph{после} SNAT. Только в этот момент известны и исходные,
и окончательные адреса, и выходной порт.

Быстрый путь важен и для производительности: полный просмотр DNAT, политики
и SNAT выполняется один раз на сеанс, а не на каждый пакет
(раздел~\ref{sec:perf}).

\subsection{Связи с другими подсистемами}
\begin{itemize}
\item \textbf{Регистрация.} Если у правила стоит флаг регистрации, после
  обработки кадра формируется запись о пакете и отправляется в подсистему
  регистрации (стрелка «записи о пакетах»). Отправка не ждёт ответа: при
  переполнении буфера запись теряется, а счётчик потерь растёт, но фильтрация
  не тормозит (ОБР-10).
\item \textbf{Контроль целостности.} Стрелка «глобальная блокировка» включает
  режим, в котором блок~1 удаляет все кадры (ЦЛС-04).
\item \textbf{Управление.} Прикладные сервисы меняют правила через API
  политики (АПИ-01\,--\,АПИ-12). Новая версия таблиц собирается целиком
  и заменяет старую одним действием (СТР-07): пакет обрабатывается либо по
  старой, либо по новой версии, но никогда по смеси.
\item \textbf{Собственные службы изделия.} SSH и HTTPS администратора приходят
  на физический порт, проходят политику и с выходом \texttt{local} попадают
  в TAP-интерфейс, а оттуда в стек ОС к sshd и web-серверу (ОБР-05). Значит,
  удалённое управление тоже защищено политикой.
\end{itemize}

% =====================================================================
\section{Подсистема управления}\label{sec:mgmt}

\subsection{Точки входа}
\begin{itemize}
\item \textbf{CLI локально}~— консоль изделия (монитор и клавиатура или
  последовательный порт). Нужна для первоначальной настройки, когда в
  политике ещё нет правил для удалённого доступа, и для аварийного режима
  (ЦЛС-05).
\item \textbf{CLI удалённо}~— та же программа CLI, запущенная через SSH
  (ДСТ-07). Шифрование канала обеспечивает SSH.
\item \textbf{Web-интерфейс}~— HTTPS, TLS 1.2 и выше (ДСТ-07): страницы
  и запросы к единому API.
\end{itemize}
Точки входа не хранят данных и не принимают решений. Каждое действие
администратора они превращают в вызов единого API и показывают ответ.

\subsection{Единый API управления}

Через этот блок проходит каждая команда, откуда бы она ни пришла.
\begin{itemize}
\item \textbf{Аутентификация и сеансы.} Вход по идентификатору и паролю
  (ДСТ-01), хранение сеансов, тайм-аут неактивности (ДСТ-06), принудительный
  сброс (УПР-12).
\item \textbf{Права ЧЗ/ТЧ.} Для каждой функции проверяется уровень прав
  (ДСТ-03). Сеансы CLI и web хранятся в одном месте, поэтому правило
  «второй администратор с ЧЗ получает ТЧ» выполняется независимо от того,
  через какой интерфейс он вошёл (ДСТ-05).
\item \textbf{Валидация.} Правила и параметры проверяются до выполнения
  (ФЛТ-05, НАТ-12, НДЖ-04). Ошибки возвращаются одинаковым текстом в CLI и web.
\item \textbf{Подтверждение опасных операций.} Выключение, перезагрузка,
  применение конфигурации, очистка журнала пакетов и удаление учётной записи
  выполняются только после подтверждения (УПР-15).
\end{itemize}

\subsection{Прикладные сервисы}

Прикладные сервисы выполняют проверенные команды и обращаются к нужной
подсистеме или данным (стрелки из блока на рисунке~\ref{fig:arch}):
\begin{itemize}
\item \textbf{правила}~— в API политики подсистемы фильтрации;
\item \textbf{конфигурация}~— сохранение, вывод и применение (УПР-09\,--\,УПР-11),
  чтение и запись в постоянные данные;
\item \textbf{журналы и состояния}~— выборки из журналов (УПР-04\,--\,УПР-07)
  через подсистему регистрации, таблица состояний (УПР-08) через фильтрацию;
\item \textbf{учётные записи}~— управление администраторами (УПР-12);
\item \textbf{сеть и ОС}~— интерфейсы, адреса, маршруты, время, перезапуск
  подсистем и изделия (УПР-13, УПР-14) через блок «Управление средой»;
\item \textbf{запрос проверки целостности}~— запуск проверки по команде
  администратора (ЦЛС-03).
\end{itemize}
Кроме того, сервисы отправляют в подсистему регистрации события о действиях
администраторов: входах, изменениях политики и параметров (РЕГ-01, РЕГ-02).

% =====================================================================
\section{Подсистема регистрации}\label{sec:log}

Подсистема принимает данные из двух источников:
\begin{itemize}
\item \textbf{события}~— от подсистемы управления и остальных подсистем
  (запуск, остановка, аварийное завершение, контроль целостности) с кодами
  из приложения~А ТЗ;
\item \textbf{записи о пакетах}~— от подсистемы фильтрации, отдельным потоком.
\end{itemize}
Она же выполняет выборки по запросам прикладных сервисов: по коду и времени
для событий, по действию, интерфейсам, правилу, адресам и портам для пакетов
(УПР-04, УПР-05). Журналы хранятся в постоянных данных. Журнал событий
средствами управления не изменяется, журнал пакетов очищается только
администратором с правами ЧЗ (РЕГ-05).

% =====================================================================
\section{Подсистема контроля целостности}\label{sec:integrity}

\begin{itemize}
\item \textbf{Что проверяется:} исполняемые файлы и конфигурационные файлы
  ПО МЭ, включая сохранённые конфигурации и базу учётных записей (ЦЛС-01).
\item \textbf{Как:} SHA-256 каждого файла сравнивается с эталоном из постоянных
  данных.
\item \textbf{Когда:} при запуске, по расписанию (от 1 до 1440 минут) и по
  запросу администратора (ЦЛС-03).
\item \textbf{Легитимные изменения.} Когда управление само меняет
  конфигурационный файл (сохранение конфигурации, изменение политики, учётных
  записей), эталон и резервная копия обновляются (ЦЛС-02).
\item \textbf{При нарушении:} событие с перечнем файлов, глобальная блокировка
  трафика и переход в аварийный режим (ЦЛС-04). В аварийном режиме доступен
  только локальный CLI с функциями восстановления из резервных копий
  (ЦЛС-05). После успешной повторной проверки блокировка снимается (ЦЛС-06).
\end{itemize}

% =====================================================================
\section{Постоянные данные}\label{sec:data}

{\small
\begin{longtable}{|L{3.2cm}|L{6.4cm}|L{5.0cm}|}
\caption{Постоянные данные}\label{tab:data}\\
\hline
\textbf{Хранилище} & \textbf{Содержимое} & \textbf{Кто пишет}\\
\hline\endfirsthead
\hline
\textbf{Хранилище} & \textbf{Содержимое} & \textbf{Кто пишет}\\
\hline\endhead
Конфигурация & Текстовый файл: правила фильтрации, DNAT, SNAT, порты,
  параметры подсистем (УПР-10). Текущая и сохранённая, а также конфигурация
  по умолчанию & Прикладные сервисы\\ \hline
Учётные записи & Идентификаторы, уровни прав, хеши паролей с солью (ДСТ-02).
  В конфигурацию не входят & Прикладные сервисы\\ \hline
Контрольные суммы & Эталонные SHA-256 и резервные копии контролируемых файлов &
  Подсистема контроля целостности\\ \hline
Два журнала & Журнал событий (неизменяемый) и журнал пакетов (очистка по ЧЗ) &
  Подсистема регистрации\\ \hline
\end{longtable}
}

% =====================================================================
\section{ОС Linux и жизненный цикл}\label{sec:os}

Блок «Управление средой» выполняет всё, что требует прав в ОС:
\begin{itemize}
\item приводит ОС в состояние из ОБР-02: у портов нет IP-адресов, они
  в режиме \texttt{promisc}, а правило nftables \texttt{netdev ingress drop}
  не пускает кадры портов в стек ядра;
\item создаёт TAP-интерфейсы и назначает на них IP-адреса изделия (ОБР-05,
  УПР-13);
\item управляет системными маршрутами и временем (УПР-14);
\item запускает подсистемы, следит за ними и перезапускает после аварийного
  завершения с регистрацией события (НДЖ-01); перезапускает подсистему или
  изделие по команде администратора.
\end{itemize}

% =====================================================================
\section{Реализация на Python 3.13}\label{sec:python}

\subsection{Процессы}

Каждая подсистема реализуется отдельной программой на Python~3.13 и запускается
как служба systemd. Подсистемы обмениваются сообщениями через UNIX-сокеты
в каталоге \texttt{/run/fw/} в формате АПИ-03/04: одна JSON-строка на запрос
или ответ.

{\small
\begin{longtable}{|L{2.8cm}|L{2.5cm}|L{5.9cm}|L{3.0cm}|}
\caption{Процессы и их соответствие схеме}\label{tab:procs}\\
\hline
\textbf{Блок схемы} & \textbf{Процесс} & \textbf{Реализация} & \textbf{Права в ОС}\\
\hline\endfirsthead
\hline
\textbf{Блок схемы} & \textbf{Процесс} & \textbf{Реализация} & \textbf{Права в ОС}\\
\hline\endhead
01 Фильтрация & \texttt{fw-filter} &
  Один поток, цикл \texttt{selectors} по сокетам портов, TAP и сокету API
  политики &
  пользователь \texttt{fwfilter}, \texttt{CAP\_NET\_RAW}, \texttt{CAP\_NET\_ADMIN} (ОБР-12)\\ \hline
02 Единый API и прикладные сервисы & \texttt{fw-mgmt} &
  \texttt{asyncio}, сервер на \texttt{/run/fw/mgmt.sock}; сеансы, права,
  валидация, вызовы остальных подсистем &
  \texttt{root} (выполняет действия «Управления средой»)\\ \hline
02 CLI локально и удалённо & \texttt{fw-cli} &
  Клиент \texttt{mgmt.sock}; на консоли~— оболочка входа, по SSH~—
  \texttt{ForceCommand} в \texttt{sshd\_config} &
  \texttt{fwcli}\\ \hline
02 Web-интерфейс & \texttt{fw-web} &
  \texttt{http.server.\allowbreak ThreadingHTTPServer} + \texttt{ssl}, страницы
  и JSON-запросы к \texttt{mgmt.sock} &
  \texttt{fwweb}, \texttt{CAP\_NET\_\allowbreak BIND\_\allowbreak SERVICE}\\ \hline
03 Регистрация & \texttt{fw-log} &
  \texttt{asyncio}; \texttt{sqlite3} (режим WAL), базы событий и пакетов &
  \texttt{fwlog}\\ \hline
04 Контроль целостности & \texttt{fw-integrity} &
  \texttt{hashlib.sha256}, таймер проверок, резервные копии &
  \texttt{root} (чтение всех файлов)\\ \hline
ОС / жизненный цикл & systemd + \texttt{fw-mgmt} &
  Юниты с \texttt{Restart=on-failure}; команды \texttt{ip}, \texttt{nft},
  \texttt{sysctl}, \texttt{ethtool}, \texttt{timedatectl} через \texttt{subprocess} &
  \texttt{root}\\ \hline
\end{longtable}
}

Перезапуск подсистем выполняет systemd, а события о нём фиксируются
так: в юнит каждой подсистемы добавляются \texttt{ExecStartPost} и
\texttt{ExecStopPost}, которые вызывают маленькую утилиту
\texttt{fw-notify}. Она передаёт в \texttt{fw-log} событие запуска (210)
или остановки (211, а при \texttt{\$SERVICE\_RESULT}~$\neq$~\texttt{success}~— 213)
с кодом завершения из переменных \texttt{\$EXIT\_CODE} и \texttt{\$EXIT\_STATUS}.

\subsection{Тракт фильтрации на Python}\label{sec:perf}

\begin{itemize}
\item \textbf{Один поток вместо многих.} Из-за GIL потоки Python не дают
  параллельной обработки, а таблица состояний должна быть общей для всех портов
  (ответ приходит на другой порт). Поэтому \texttt{fw-filter} обрабатывает все
  порты в одном цикле \texttt{selectors} (epoll). Блокировки не нужны,
  и требование безопасного изменения таблиц (ОБР-12) выполняется
  автоматически.
\item \textbf{Атомарная замена политики} (СТР-07). Команды API обрабатываются
  в том же цикле между кадрами. Новая версия таблиц собирается в отдельный
  объект, а затем одним присваиванием подменяет ссылку на текущую версию.
\item \textbf{Сокеты портов.}
  \texttt{socket.socket(AF\_PACKET,\allowbreak{} SOCK\_RAW,\allowbreak{} htons(ETH\_P\_ALL))} с привязкой
  к интерфейсу. Обязательно устанавливается опция \texttt{PACKET\_\allowbreak IGNORE\_\allowbreak OUTGOING}
  (значение 23, уровень \texttt{SOL\_PACKET}~= 263). Иначе сокет получает
  обратно кадры, которые сам же отправил, и они пойдут по кругу.
\item \textbf{TAP.} Открывается \texttt{/dev/net/tun} и вызывается
  \texttt{fcntl.ioctl} с \texttt{TUNSETIFF} (флаги \texttt{IFF\_TAP | IFF\_NO\_PI}).
  Дескриптор добавляется в тот же цикл \texttt{selectors}.
\item \textbf{Разбор без копирования.} Кадр читается в заранее выделенный
  буфер (\texttt{recv\_into}), заголовки разбираются через \texttt{memoryview}
  и заранее скомпилированные \texttt{struct.Struct}.
\item \textbf{Правила в «скомпилированном» виде.} При загрузке политики
  адреса переводятся в целые числа, подсети~— в пары (адрес, маска), порты~—
  в диапазоны. Проверка правила сводится к сравнению целых чисел.
\item \textbf{Таблица состояний}~— словарь с ключом (протокол, адреса, порты),
  по одной записи на каждое направление.
\item \textbf{Контрольные суммы после NAT} пересчитываются инкрементально
  (RFC~1624): меняются только изменённые 16-битные слова, а не весь пакет.
\end{itemize}

\textbf{Оценка производительности (ОБР-13).} 100~Мбит/с при кадрах 1518 байт~—
это примерно 8\,100 кадров в секунду (плюс около 4\,000 коротких подтверждений
TCP в обратную сторону). На один кадр приходится около 80~мкс. Пакеты
установленного сеанса идут по быстрому пути (поиск в словаре, трансляция,
отправка), и это укладывается в бюджет даже на одном ядре. Полный просмотр
100 правил DNAT, политики и SNAT выполняется только для первого пакета сеанса.
Поэтому для теста \texttt{iperf3} правило должно быть с контролем состояния.
Оценку нужно подтвердить замером на стенде.

\subsection{Используемые модули}

Используется только стандартная библиотека Python, поэтому дополнительные
пакеты устанавливать не нужно.

{\small
\begin{longtable}{|L{4.4cm}|L{10.2cm}|}
\caption{Модули стандартной библиотеки}\label{tab:modules}\\
\hline
\textbf{Модуль} & \textbf{Для чего}\\
\hline\endfirsthead
\hline
\textbf{Модуль} & \textbf{Для чего}\\
\hline\endhead
\texttt{socket}, \texttt{selectors} & \texttt{AF\_PACKET}, UNIX-сокеты, цикл обработки кадров\\ \hline
\texttt{fcntl}, \texttt{os} & создание TAP-интерфейсов, чтение и запись дескриптора\\ \hline
\texttt{struct}, \texttt{ipaddress} & разбор заголовков, адреса и подсети\\ \hline
\texttt{asyncio} & серверы \texttt{fw-mgmt} и \texttt{fw-log}\\ \hline
\texttt{json} & протокол обмена и формат конфигурации\\ \hline
\texttt{hashlib} & SHA-256 (ЦЛС-01); \texttt{hashlib.scrypt} для хешей паролей с солью (ДСТ-02)\\ \hline
\texttt{secrets}, \texttt{hmac} & токены сеансов, сравнение без утечки по времени\\ \hline
\texttt{ssl}, \texttt{http.server} & HTTPS-сервер, \texttt{minimum\_version = TLSv1\_2} (ДСТ-07)\\ \hline
\texttt{sqlite3} & журналы событий и пакетов, выборки по критериям\\ \hline
\texttt{subprocess} & вызов \texttt{ip}, \texttt{nft}, \texttt{sysctl}, \texttt{ethtool}, \texttt{timedatectl}, \texttt{systemctl}\\ \hline
\texttt{logging}, \texttt{argparse} & отладочный вывод и разбор аргументов CLI\\ \hline
\end{longtable}
}

% =====================================================================

\end{document}
